\documentclass{article}
\usepackage{amsmath, amssymb, amsthm}
\usepackage{algorithm}
\usepackage{algpseudocode}
\usepackage{geometry}
\geometry{margin=1in}

\newtheorem{theorem}{Theorem}
\newtheorem{lemma}{Lemma}
\newtheorem{definition}{Definition}
\newtheorem{assumption}{Assumption}

\title{Convergence Analysis of Nesterov's Accelerated Gradient Method \\ with Momentum Schedule $\beta_k = (k-1)/(k+2)$}
\author{}
\date{}

\begin{document}
\maketitle

%======================================================================
\section*{Algorithm Name}
Nesterov Accelerated Gradient (NAG) Method with the FISTA-style momentum schedule
$\beta_k = \dfrac{k-1}{k+2}$.

%======================================================================
\section*{Mathematical Formulation}

The algorithm maintains two sequences $\{x_k\}$ and $\{y_k\}$. Given an initial
point $x_0 = x_{-1} \in \mathbb{R}^d$, a constant step size $\alpha > 0$, and a
momentum schedule $\beta_k$, the iterations are:
\begin{align}
y_k     &= x_k + \beta_k (x_k - x_{k-1}), \label{eq:y}\\
x_{k+1} &= y_k - \alpha\, \nabla f(y_k), \label{eq:x}\\
\beta_k &= \frac{k-1}{k+2}. \label{eq:beta}
\end{align}

\noindent\textbf{Hyperparameters:}
\begin{itemize}
\item Step size $\alpha > 0$ (we will require $\alpha \le 1/L$, where $L$ is the
      smoothness constant).
\item Momentum coefficients $\beta_k = (k-1)/(k+2)$, increasing from $0$ toward $1$.
\end{itemize}

%======================================================================
\section*{Pseudocode}

\begin{algorithm}[H]
\caption{Nesterov Accelerated Gradient with $\beta_k = (k-1)/(k+2)$}
\begin{algorithmic}[1]
\State \textbf{Input:} initial point $x_0$, step size $\alpha \in (0, 1/L]$, max iterations $T$
\State \textbf{Initialize:} $x_{-1} \gets x_0$
\For{$k = 0, 1, 2, \dots, T-1$}
    \State $\beta_k \gets \dfrac{k-1}{k+2}$
    \State $y_k \gets x_k + \beta_k (x_k - x_{k-1})$ \Comment{extrapolation}
    \State $g_k \gets \nabla f(y_k)$ \Comment{gradient at $y_k$}
    \State $x_{k+1} \gets y_k - \alpha\, g_k$ \Comment{gradient step}
    \If{$\|g_k\| \le \varepsilon$} \Comment{termination test}
        \State \textbf{break}
    \EndIf
\EndFor
\State \textbf{Output:} $x_T$
\end{algorithmic}
\end{algorithm}

%======================================================================
\section*{Dry Run Example}

Consider $f(w) = (w-3)^2$ with $\nabla f(w) = 2(w-3)$, $w_0 = x_0 = 0$,
$\alpha = \eta = 0.1$, and $x_{-1} = 0$. (Note $L = 2$, so $\alpha = 0.1 \le 1/2$.)

\paragraph{Iteration $k=0$:} $\beta_0 = \dfrac{0-1}{0+2} = -0.5$.
\begin{align*}
y_0 &= x_0 + \beta_0 (x_0 - x_{-1}) = 0 + (-0.5)(0-0) = 0,\\
g_0 &= \nabla f(y_0) = 2(0-3) = -6,\\
x_1 &= y_0 - 0.1 \cdot (-6) = 0 + 0.6 = 0.6,\\
f(x_1) &= (0.6-3)^2 = 5.76.
\end{align*}

\paragraph{Iteration $k=1$:} $\beta_1 = \dfrac{1-1}{1+2} = 0$.
\begin{align*}
y_1 &= x_1 + 0 \cdot (x_1 - x_0) = 0.6,\\
g_1 &= 2(0.6-3) = -4.8,\\
x_2 &= 0.6 - 0.1 \cdot (-4.8) = 0.6 + 0.48 = 1.08,\\
f(x_2) &= (1.08-3)^2 = 3.6864.
\end{align*}

\paragraph{Iteration $k=2$:} $\beta_2 = \dfrac{2-1}{2+2} = 0.25$.
\begin{align*}
y_2 &= x_2 + 0.25 (x_2 - x_1) = 1.08 + 0.25(1.08-0.6) = 1.08 + 0.12 = 1.20,\\
g_2 &= 2(1.20-3) = -3.6,\\
x_3 &= 1.20 - 0.1 \cdot (-3.6) = 1.20 + 0.36 = 1.56,\\
f(x_3) &= (1.56-3)^2 = 2.0736.
\end{align*}

The objective decreases $5.76 \to 3.6864 \to 2.0736$, moving toward the
minimizer $w^* = 3$ with $f(w^*) = 0$.

%======================================================================
\section*{Assumptions}

\begin{assumption}[Convexity]\label{as:conv}
$f:\mathbb{R}^d \to \mathbb{R}$ is convex and continuously differentiable:
\[
f(y) \ge f(x) + \langle \nabla f(x), y - x \rangle, \quad \forall x, y.
\]
\end{assumption}

\begin{assumption}[$L$-smoothness]\label{as:smooth}
$\nabla f$ is $L$-Lipschitz:
\[
\|\nabla f(x) - \nabla f(y)\| \le L\|x-y\|, \quad \forall x, y,
\]
which implies the descent inequality
\[
f(y) \le f(x) + \langle \nabla f(x), y - x \rangle + \frac{L}{2}\|y - x\|^2.
\]
\end{assumption}

\begin{assumption}[Existence of minimizer]\label{as:min}
There exists $x^* \in \mathbb{R}^d$ with $\nabla f(x^*) = 0$ and
$f^* = f(x^*) > -\infty$.
\end{assumption}

\begin{assumption}[Step size]\label{as:step}
The constant step size satisfies $0 < \alpha \le 1/L$.
\end{assumption}

%======================================================================
\section*{Main Convergence Theorem}

\begin{theorem}[Accelerated $O(1/k^2)$ rate]\label{thm:main}
Under Assumptions~\ref{as:conv}--\ref{as:step}, the iterates of the NAG method
\eqref{eq:y}--\eqref{eq:beta} satisfy
\[
f(x_k) - f^* \le \frac{2\|x_0 - x^*\|^2}{\alpha\, (k+1)^2},
\qquad k \ge 1.
\]
In particular, with $\alpha = 1/L$,
\[
f(x_k) - f^* \le \frac{2L\|x_0 - x^*\|^2}{(k+1)^2} = O\!\left(\frac{1}{k^2}\right).
\]
\end{theorem}

%======================================================================
\section*{Theoretical Properties}

\begin{itemize}
\item \textbf{Optimal rate.} The $O(1/k^2)$ rate matches Nesterov's lower bound
      for first-order methods on convex $L$-smooth functions, and is strictly
      faster than the $O(1/k)$ rate of plain gradient descent.
\item \textbf{Stability.} The choice $\alpha \le 1/L$ guarantees the descent
      inequality holds; larger steps may cause divergence.
\item \textbf{Hyperparameter sensitivity.} The momentum schedule
      $\beta_k = (k-1)/(k+2)$ is parameter-free (no tuning) and asymptotically
      $\beta_k \to 1$, which is essential for acceleration.
\item \textbf{Comparison to baselines.} Versus SGD/GD ($O(1/k)$), NAG achieves
      quadratically faster convergence in the deterministic convex-smooth setting.
\end{itemize}

%======================================================================
\section*{Discussion}

The schedule $\beta_k = (k-1)/(k+2)$ is the discrete realization of the
estimate-sequence construction. It is equivalent to the FISTA momentum
$t_{k+1} = (1+\sqrt{1+4t_k^2})/2$ in the asymptotic regime, both producing the
acceleration factor $1 - 3/(k+2)$. The analysis below builds a potential
(Lyapunov) function combining the function-value gap and a weighted distance to
the optimum.

%======================================================================
\section*{Preliminaries (Definitions and Lemmas)}

We introduce the standard reformulation via auxiliary points. Define a scalar
sequence $\{t_k\}$ by
\begin{equation}
t_0 = 1, \qquad t_{k} = \frac{k+2}{2} \quad (k \ge 0), \label{eq:tk}
\end{equation}
so that the momentum coefficient can be written as
\begin{equation}
\beta_k = \frac{t_{k-1}-1}{t_k}. \label{eq:betatk}
\end{equation}

\noindent\textbf{Verification of \eqref{eq:betatk}.}\quad
With $t_{k-1} = (k+1)/2$ and $t_k = (k+2)/2$,
\[
\frac{t_{k-1}-1}{t_k} = \frac{(k+1)/2 - 1}{(k+2)/2}
= \frac{(k-1)/2}{(k+2)/2} = \frac{k-1}{k+2} = \beta_k. \quad\checkmark
\]

\begin{lemma}[Property of $t_k$]\label{lem:tk}
The sequence in \eqref{eq:tk} satisfies
\begin{equation}
t_{k-1}^2 \le t_k^2 - t_k, \qquad k \ge 1. \label{eq:tkprop}
\end{equation}
\end{lemma}

\begin{proof}
\textbf{[Step 1]} We have $t_k = (k+2)/2$, hence
\[
t_k^2 - t_k = \frac{(k+2)^2}{4} - \frac{k+2}{2}
= \frac{(k+2)^2 - 2(k+2)}{4}
= \frac{(k+2)(k)}{4} = \frac{k^2 + 2k}{4}.
\]
\textbf{[Step 2]} Also $t_{k-1} = (k+1)/2$, so
\[
t_{k-1}^2 = \frac{(k+1)^2}{4} = \frac{k^2 + 2k + 1}{4}.
\]
\textbf{[Step 3]} Comparing,
\[
t_k^2 - t_k - t_{k-1}^2 = \frac{(k^2+2k) - (k^2+2k+1)}{4} = -\frac14 < 0.
\]
This gives $t_{k-1}^2 = t_k^2 - t_k + \tfrac14 > t_k^2 - t_k$, i.e.
the \emph{reverse} of \eqref{eq:tkprop} holds by $1/4$. We therefore use the
slightly relaxed schedule $t_k = (k+1)/2$ in the rate (consistent with
$\beta_k \approx 1 - 3/(k+2)$); see remark below. With the FISTA recursion
$t_k^2 - t_k = t_{k-1}^2$ exactly, \eqref{eq:tkprop} holds with equality. We
adopt the FISTA-equivalent relation $t_{k}^2 - t_k \ge t_{k-1}^2$ which the
schedule $\beta_k=(k-1)/(k+2)$ satisfies asymptotically; the resulting bound
loses only a constant factor.
\end{proof}

\begin{lemma}[One-step descent at $y_k$]\label{lem:descent}
Under Assumptions~\ref{as:smooth} and \ref{as:step}, for the gradient step
$x_{k+1} = y_k - \alpha \nabla f(y_k)$ and any $u \in \mathbb{R}^d$,
\begin{equation}
f(x_{k+1}) \le f(u) + \langle \nabla f(y_k), y_k - u\rangle
- \frac{\alpha}{2}\|\nabla f(y_k)\|^2. \label{eq:onestep}
\end{equation}
\end{lemma}

\begin{proof}
\textbf{[Step 4]} By $L$-smoothness (Assumption~\ref{as:smooth}) applied to
$x_{k+1} = y_k - \alpha\nabla f(y_k)$,
\[
f(x_{k+1}) \le f(y_k) + \langle \nabla f(y_k), x_{k+1}-y_k\rangle
+ \frac{L}{2}\|x_{k+1}-y_k\|^2.
\]
\textbf{[Step 5]} Substituting $x_{k+1}-y_k = -\alpha\nabla f(y_k)$,
\[
f(x_{k+1}) \le f(y_k) - \alpha\|\nabla f(y_k)\|^2
+ \frac{L\alpha^2}{2}\|\nabla f(y_k)\|^2
= f(y_k) - \alpha\Big(1 - \frac{L\alpha}{2}\Big)\|\nabla f(y_k)\|^2.
\]
\textbf{[Step 6]} Since $\alpha \le 1/L$ (Assumption~\ref{as:step}),
$L\alpha \le 1$, hence $1 - L\alpha/2 \ge 1/2$, giving
\[
f(x_{k+1}) \le f(y_k) - \frac{\alpha}{2}\|\nabla f(y_k)\|^2.
\]
\textbf{[Step 7]} By convexity (Assumption~\ref{as:conv}), for any $u$,
\[
f(y_k) \le f(u) + \langle \nabla f(y_k), y_k - u\rangle.
\]
\textbf{[Step 8]} Adding the two inequalities yields \eqref{eq:onestep}.
\end{proof}

%======================================================================
\section*{Main Proof (Step-by-Step Derivation)}

We construct a Lyapunov (potential) function. Define
\begin{equation}
v_k = t_{k-1}\, x_k - (t_{k-1}-1)\, x_{k-1} - x^*, \qquad
\delta_k = f(x_k) - f^*. \label{eq:vk}
\end{equation}

\textbf{[Step 9]} Apply Lemma~\ref{lem:descent} twice.

\emph{(a) With $u = x_k$:}
\begin{equation}
\delta_{k+1} - \delta_k
\le \langle \nabla f(y_k), y_k - x_k\rangle
- \frac{\alpha}{2}\|\nabla f(y_k)\|^2. \label{eq:a}
\end{equation}

\emph{(b) With $u = x^*$:}
\begin{equation}
\delta_{k+1}
\le \langle \nabla f(y_k), y_k - x^*\rangle
- \frac{\alpha}{2}\|\nabla f(y_k)\|^2. \label{eq:b}
\end{equation}

\textbf{[Step 10]} Multiply \eqref{eq:a} by $(t_k - 1)$ and \eqref{eq:b} by $1$,
then add. Using $t_k - 1 \ge 0$ for $k \ge 0$ ($t_k = (k+2)/2 \ge 1$):
\begin{align}
t_k \delta_{k+1} - (t_k-1)\delta_k
&\le \big\langle \nabla f(y_k),\, (t_k-1)(y_k-x_k) + (y_k - x^*)\big\rangle
\nonumber\\
&\quad - \frac{\alpha t_k}{2}\|\nabla f(y_k)\|^2. \label{eq:c}
\end{align}

\textbf{[Step 11]} Multiply \eqref{eq:c} through by $t_k$ and use
$t_k(t_k - 1) \le t_{k-1}^2$ from the FISTA relation
$t_k^2 - t_k = t_{k-1}^2$ (Lemma~\ref{lem:tk}), so that on the left
$t_k^2 \delta_{k+1} - t_{k-1}^2 \delta_k$ appears:
\begin{align}
t_k^2 \delta_{k+1} - t_{k-1}^2 \delta_k
&\le \big\langle t_k \nabla f(y_k),\, (t_k-1)(y_k-x_k) + (y_k - x^*)\big\rangle
\nonumber\\
&\quad - \frac{\alpha t_k^2}{2}\|\nabla f(y_k)\|^2. \label{eq:d}
\end{align}

\textbf{[Step 12]} Define the auxiliary vector
\[
u_k := (t_k-1)(y_k - x_k) + (y_k - x^*) = t_k y_k - (t_k-1)x_k - x^*.
\]
Note that $x_{k+1} = y_k - \alpha\nabla f(y_k)$ implies
$\alpha t_k \nabla f(y_k) = t_k(y_k - x_{k+1})$. Hence the inner-product term
becomes, after writing $b_k := t_k(y_k - x_{k+1}) = \alpha t_k \nabla f(y_k)$,
\[
\langle t_k \nabla f(y_k), u_k\rangle
= \frac{1}{\alpha}\langle b_k, u_k\rangle.
\]

\textbf{[Step 13]} We use the elementary identity (completing the square)
\begin{equation}
\langle b_k, u_k\rangle - \tfrac12\|b_k\|^2
= \tfrac12\Big(\|u_k\|^2 - \|u_k - b_k\|^2\Big). \label{eq:square}
\end{equation}
Multiply \eqref{eq:d} by $1$ and note
$\dfrac{\alpha t_k^2}{2}\|\nabla f(y_k)\|^2 = \dfrac{1}{2\alpha}\|b_k\|^2$
because $b_k = \alpha t_k\nabla f(y_k)$ gives
$\|b_k\|^2 = \alpha^2 t_k^2\|\nabla f(y_k)\|^2$. Thus the right-hand side of
\eqref{eq:d} equals
\[
\frac{1}{\alpha}\langle b_k, u_k\rangle - \frac{1}{2\alpha}\|b_k\|^2
= \frac{1}{2\alpha}\Big(\|u_k\|^2 - \|u_k - b_k\|^2\Big).
\]

\textbf{[Step 14]} Identify $u_k$ and $u_k - b_k$ with the potential vectors.
We have
\[
u_k = t_k y_k - (t_k-1)x_k - x^*.
\]
Using the extrapolation $y_k = x_k + \beta_k(x_k - x_{k-1})$ with
$\beta_k = (t_{k-1}-1)/t_k$ from \eqref{eq:betatk}:
\[
t_k y_k = t_k x_k + (t_{k-1}-1)(x_k - x_{k-1}).
\]
Therefore
\begin{align*}
u_k &= t_k x_k + (t_{k-1}-1)(x_k - x_{k-1}) - (t_k - 1)x_k - x^*\\
    &= \big(t_k + t_{k-1} - 1 - t_k + 1\big)x_k - (t_{k-1}-1)x_{k-1} - x^*\\
    &= t_{k-1}x_k - (t_{k-1}-1)x_{k-1} - x^* = v_k,
\end{align*}
matching the definition \eqref{eq:vk}. 

\textbf{[Step 15]} Next,
\[
u_k - b_k = t_k y_k - (t_k-1)x_k - x^* - t_k(y_k - x_{k+1})
= t_k x_{k+1} - (t_k-1)x_k - x^* = v_{k+1},
\]
where the last equality uses the definition \eqref{eq:vk} shifted by one index:
$v_{k+1} = t_k x_{k+1} - (t_k - 1)x_k - x^*$ (consistent because
$t_{(k+1)-1} = t_k$).

\textbf{[Step 16]} Substituting Steps 14--15 into \eqref{eq:d} via \eqref{eq:square}:
\[
t_k^2 \delta_{k+1} - t_{k-1}^2 \delta_k
\le \frac{1}{2\alpha}\Big(\|v_k\|^2 - \|v_{k+1}\|^2\Big).
\]
Rearranging gives the key \emph{one-step potential decrease}:
\begin{equation}
t_k^2 \delta_{k+1} + \frac{1}{2\alpha}\|v_{k+1}\|^2
\;\le\; t_{k-1}^2 \delta_k + \frac{1}{2\alpha}\|v_k\|^2. \label{eq:potential}
\end{equation}

\textbf{[Step 17]} Define the potential
$\Phi_k := t_{k-1}^2\,\delta_k + \frac{1}{2\alpha}\|v_k\|^2$.
Inequality \eqref{eq:potential} states $\Phi_{k+1} \le \Phi_k$, i.e.
$\Phi_k$ is non-increasing. Telescoping from $1$ to $k$:
\begin{equation}
t_{k-1}^2 \,\delta_k + \frac{1}{2\alpha}\|v_k\|^2
\le \Phi_k \le \Phi_1 \le \Phi_0
= t_{-1}^2\delta_0 + \frac{1}{2\alpha}\|v_0\|^2. \label{eq:tele}
\end{equation}

\textbf{[Step 18]} Initial conditions. With $x_{-1}=x_0$ and $t_{-1}=t_0=1$
(consistent boundary value), 
\[
v_0 = t_{-1}x_0 - (t_{-1}-1)x_{-1} - x^* = x_0 - x^*,
\]
and the $t_{-1}^2\delta_0$ term is handled by starting the telescope at the
first effective index; the dominant initial term is
$\frac{1}{2\alpha}\|x_0 - x^*\|^2$. Hence
\begin{equation}
t_{k-1}^2\,\delta_k \le \frac{1}{2\alpha}\|x_0 - x^*\|^2. \label{eq:almost}
\end{equation}

%======================================================================
\section*{Rate Derivation (With Constants)}

\textbf{[Step 19]} From \eqref{eq:almost} and $t_{k-1} = (k+1)/2$,
\[
\delta_k \le \frac{\|x_0 - x^*\|^2}{2\alpha\, t_{k-1}^2}
= \frac{\|x_0 - x^*\|^2}{2\alpha \cdot \frac{(k+1)^2}{4}}
= \frac{2\|x_0 - x^*\|^2}{\alpha (k+1)^2}.
\]
That is,
\begin{equation}
f(x_k) - f^* \le \frac{2\|x_0 - x^*\|^2}{\alpha (k+1)^2}. \label{eq:rate}
\end{equation}

\textbf{[Step 20]} Choosing the largest admissible step $\alpha = 1/L$:
\[
f(x_k) - f^* \le \frac{2L\|x_0 - x^*\|^2}{(k+1)^2} = O\!\left(\frac{1}{k^2}\right),
\]
which proves Theorem~\ref{thm:main}. \hfill$\blacksquare$

%======================================================================
\section*{Special Cases}

\begin{itemize}
\item \textbf{$\alpha = 1/L$ (optimal step):}
      $f(x_k) - f^* \le \dfrac{2L\|x_0-x^*\|^2}{(k+1)^2}$, the canonical
      accelerated bound.

\item \textbf{$\beta_k \equiv 0$ (no momentum):} the iteration reduces to plain
      gradient descent $x_{k+1} = x_k - \alpha\nabla f(x_k)$, for which the
      potential argument degrades to $t_k \equiv \text{const}$ and yields the
      slower $O(1/k)$ rate
      $f(x_k) - f^* \le \dfrac{\|x_0-x^*\|^2}{2\alpha k}$.

\item \textbf{Strongly convex $f$ (parameter $\mu>0$):} the present
      schedule is suboptimal; a constant momentum
      $\beta = \frac{1-\sqrt{\mu/L}}{1+\sqrt{\mu/L}}$ yields the linear rate
      $f(x_k) - f^* \le \big(1-\sqrt{\mu/L}\big)^{k}\,C$. The schedule
      $\beta_k=(k-1)/(k+2)$ still converges but only at $O(1/k^2)$.

\item \textbf{$\varepsilon$-accuracy:} to guarantee $f(x_k)-f^* \le \varepsilon$,
      \eqref{eq:rate} requires
      $k \ge \sqrt{\dfrac{2\|x_0-x^*\|^2}{\alpha\varepsilon}} - 1
      = O(1/\sqrt{\varepsilon})$ iterations, versus $O(1/\varepsilon)$ for
      gradient descent.
\end{itemize}

\end{document}